\scanpage{16}
\unsure{General non-sens} sur Frobenius

1) Soit
\diagram{\begin{tikzcd}[ampersand replacement=\&]
X\arrow[r,"f_X"]\arrow[d,"g"']\&X\arrow[d,"g"]\\
Y\arrow[r,"f_Y"']\&Y
\end{tikzcd}}
un diagramme commutatif, à isom donné près, de topos. Si $E$ est un \deleted{complexe de} faisceaux \deleted{abéliens} \inserted{d'ens} sur $X$, on a \deleted{(sous des conditions, en degré \unsure{convenable})} un hom de changement de base
\[
\text{(1.1)}\qquad c_E:f_Y^*(g_*(E))\longrightarrow g_*(f_X^*(E))
\]
Si donc on a encore un hom
\[
\text{(1.2)}\qquad F_E:f_X^*(E)\unsure{-|\longrightarrow}E
\]
on en déduit
\[
g_*(F_E):g_*f_X^*E\unsure{-|\longrightarrow}g_*(E)
\]
d'où un hom composé
\[
\text{(1.3)}\qquad\deleted{\unsure{c_E}}\quad g_*(F_E)\circ c_E:f_Y^*(g_*(E))\longrightarrow g_*(E)
\]
\deleted{En d'autres termes, on le \unsure{fait}}

Si $E$ faisceau de groupes, et (1.2) hom de tels, alors (1.3) aussi, \inserted{car (1.1) l'est.} Fait analogue pour $E$ dans une catégorie dérivée.

De $\overline{\mathrm m}$, considérons un faisceau $G$ sur $Y$, d'où un isom
\[
\text{(1.4)}\qquad c^G:f_X^*(g^*(G))\xrightarrow{\sim}g^*(f_Y^*(G))
\]
d'où, si on se donne un hom
\[
\text{(1.5)}\qquad F_G:f_Y^*(G)\longrightarrow G
\]
\margin{\uline{d'où}}
\mspage{15}
